\documentclass[11pt]{article}
\usepackage[margin=1in]{geometry}
\usepackage{booktabs}
\usepackage{longtable}
\usepackage{amsmath}
\usepackage{hyperref}
\hypersetup{colorlinks=true, linkcolor=blue, citecolor=blue, urlcolor=blue}

\title{Can the Minutes Pipeline Go Region-Wide?\\
\large The Bay Area reconnaissance sprint, June 2026 --- consolidated findings from nine probes}
\author{Dan Post}
\date{August 2026}

\begin{document}
\maketitle

\begin{center}
\fbox{\begin{minipage}{0.94\textwidth}
\small \textbf{Standing document.} The single memo for \texttt{output/bay\_area\_recon/}. It
replaces the nine per-probe memos that used to sit one per directory; each probe directory keeps
its CSVs, scripts and \texttt{.md} reports, and \S6 says which file came from where. Where a
narrative \texttt{.md} report disagrees with its own CSV, the CSV is authoritative and the
disagreement is recorded in \S7 rather than quietly resolved.
\end{minipage}}
\end{center}

\section{The question, the answer, and what is still blocked}

The project studies the \emph{market for housing regulation}: an item-level dataset of
discretionary land-use decisions built from San Francisco Planning Commission minutes
(1998--present), read through a model with two margins --- the \emph{by-right envelope}
$\bar z_j$ a locality allows without a hearing, and the \emph{discretionary} margin $\tau_j$
applied to everything else. The headline test is substitution: when the state forces the by-right
margin open, do localities claw restriction back through discretion? That test needs neighbours,
so this sprint asked one logistical question --- \textbf{can the SF pipeline be extended to the
roughly 109 land-use regulators of the nine-county ABAG Bay Area?} Nine probes ran in June 2026.

\paragraph{Yes, with one hard qualification and one unmade decision.} Retrieval is tractable: all
109 localities were surveyed and keyed to a Census GEOID, the region's meeting portals collapse
to four vendor families covering 87 of the 109, and only \textbf{three} sites region-wide are
genuinely bot-blocked. Depth is adequate for the window that matters: \textbf{36 localities} have
minutes reaching 2016 or earlier, straddling the 2017--2023 preemption ramp, and 36 is a floor,
not a ceiling. The spatial side is better --- \textbf{67 of 109} publish a downloadable zoning
GIS layer and a statewide layer backstops the rest, so address-to-district matching is a script,
not an RA-hours problem. The treatment variable is clean: all 25 estimation-sample localities
carry sourced builder's-remedy exposure dates spanning \textbf{0.0 to 38.4 months}.

The qualification is extraction. The SF heuristics \textbf{do not transfer}: run unmodified on a
real Daly City minutes document, \texttt{autoextract.extract()} produced two correct fields out
of roughly 35, and three of the wrong ones were populated with plausible-looking garbage.
Region-wide expansion is a research programme with a per-jurisdiction (or LLM) extractor and a
human in the loop, not a scripting exercise.

The unmade decision gates the whole right-hand side: \textbf{does the by-right envelope mean
``principally permitted use'' or ``no discretionary review of any kind''?} Fremont and San Mateo
both apply design or site review to uses their codes mark permitted, so the code's ``P'' column
is strictly wider than the model's ministerial envelope. Nothing downstream is well defined until
this is settled, and no further data collection will settle it.

\paragraph{Blocked, in order of consequence.} (i) The envelope definition above. (ii)
\textbf{San Jose}, whose 2005--2018 minutes are confirmed to exist but sit behind Akamai --- a
budget call, not a measurement. (iii) The \textbf{41 localities with unknown archive depth}, of
which 12 are CivicPlus sites needing one browser session to read a committee ID. (iv)
\textbf{Sample width}: both the HCD panel and the pre-period probe were fixed at 25 localities,
and depth work has since raised confirmed $\leq$2016 breadth to 36. (v) The \textbf{pre-period
envelope for 10 of 25} sample localities, a per-publisher superseded-version dig rather than
archival research.

\section{The frame: 109 land-use regulators, keyed and classified}

Everything joins to \texttt{bay\_area\_census/bay\_area\_locality\_census.csv} --- 109 rows
$\times$ 20 columns, one per entity that actually regulates land use: 101 incorporated cities and
towns, the eight county governments, and San Francisco as a consolidated reference row (85 typed
\texttt{city}, 16 \texttt{town}, 8 \texttt{county}). The city frame was supplied rather than
agent-generated, so completeness is not in question; eight per-county shallow surveys in
\texttt{raw/} (Alameda 15, Contra Costa 20, Marin 12, Napa 6, San Mateo 21, Santa Clara 16,
Solano 8, Sonoma 10 $=$ 108 localities) are the provenance for every cell. Each row records the
regulating body, minutes URL, vendor platform, code host, an access status per side, and --- the
point of the exercise --- a \emph{solution class}: the single engineering effort that would
unlock it, which turns a 109-item to-do list into a handful of adapters.

\paragraph{The region collapses into a few adapters.} Minutes side:
\texttt{granicus\_family\_clean} 41, \texttt{civicplus\_clean} 26, \texttt{custom\_cms} 10,
\texttt{civicclerk\_modern\_api} 10, \texttt{primegov\_portal} 4, \texttt{civicweb\_portal} 3,
\texttt{unknown} 3, \texttt{civicplus\_akamai} 2, \texttt{escribe\_portal} 2,
\texttt{onbase\_portal} 2, \texttt{municode\_meetings\_portal} 2, and one each of
\texttt{novusagenda\_portal}, \texttt{laserfiche\_portal}, \texttt{granicus\_family\_akamai} and
\texttt{reference\_done} (San Francisco). Four classes --- Granicus-family, CivicPlus, CivicClerk
and the custom-CMS tail --- cover \textbf{87 of 109}, and a single Granicus/Legistar adapter
unlocks 41 on its own: the highest-ROI engineering in the project. Zoning code hosts cluster
harder: \texttt{municode\_api} 44, \texttt{codepublishing\_clean} 23, \texttt{ecode360\_clean}
14, \texttt{publiclaw\_clean} 8, \texttt{amlegal\_or\_ecode\_403} 7, \texttt{unknown} 6,
\texttt{city\_site} 3, \texttt{qcode\_clean} 2, and one each \texttt{encodeplus} and
\texttt{reference\_done} --- three hosts cover 81 of 109.

\paragraph{Access is largely solved, and the Akamai scare was a myth.} Measured
\texttt{minutes\_access}: \texttt{clean} 78, \texttt{js\_shell} 16, \texttt{unknown} 10,
\texttt{akamai\_403} 3, \texttt{unsupported\_platform} 2. The 16 \texttt{js\_shell} rows are not
blocked --- their content sits behind a modern API (CivicClerk, PrimeGov, IQM2, OnBase, Municode
Meetings) and needs an API client. Zoning access: \texttt{clean} 94, \texttt{unknown} 14,
\texttt{akamai\_403} 1. Tractability tier: \texttt{needs\_adapter} 69, \texttt{clean} 35,
\texttt{access\_blocked} 3, \texttt{unknown} 2. The census was written on the premise that the
roughly 27 CivicPlus localities might \emph{all} sit behind Akamai --- a class-wide wall. A
header-only sweep of all \textbf{28} CivicPlus hosts with an ordinary browser User-Agent (no
evasion, no header tricks, no IP rotation) found \textbf{25 clean}, \textbf{2 Akamai 403s}
(Fremont, Portola Valley) and \textbf{1 stale-URL 404} (Antioch, behind Cloudflare --- a wrong
path, not a bot wall). Akamai is a per-site subscription, not a platform property; counting San
Jose, the region has \textbf{three} genuinely blocked sites. The working retrieval path was
demonstrated end-to-end: fetch \texttt{/AgendaCenter} with a normal User-Agent, parse the
embedded \texttt{ViewFile/Minutes/\_\dots} links, walk \texttt{/AgendaCenter/PreviousVersions}
for history. The 19\,KB Saratoga minutes PDF in \texttt{bay\_area\_census/} proves the path
yields a real document.

\paragraph{The join key is solid; the body names are not.} All 109 rows carry a real GEOID and
zero are \texttt{to\_verify}; names were never the join, because name joins fail silently here
(NZLUD stores Redwood City as ``Redwood''). Row confidence: \texttt{med} 64, \texttt{high} 41,
\texttt{low} 4. Minutes are recorded as posted for 104 localities, \texttt{unknown} for 4,
\texttt{sparse} for 1; none was confirmed paper-only, so there is no digitization emergency, and
the five whose platform could not be identified are Colma, Monte Sereno, Newark, Fairfax and San
Rafael. The regulating body is often \emph{not} a ``Planning Commission'' --- Planning Board
(Alameda), Planning \& Zoning Commission (Albany), Environmental Planning Commission (Mountain
View), Planning \& Transportation Commission (Palo Alto, San Carlos), a combined Planning
Commission and Board of Zoning Adjustments (San Leandro), Permit Sonoma --- so cross-jurisdiction
matching on body name will break.

\paragraph{Left open here.} Ten non-CivicPlus \texttt{minutes\_access = unknown} rows were never
probed; six zoning hosts and three minutes platforms remain \texttt{unknown}; eight localities
run dual or migrating systems where the authoritative record is not established (Foster City, San
Pablo, Menlo Park, El Cerrito, Pinole, Union City, Brentwood, Santa Clara County); Antioch's
correct agenda URL is still wrong in the CSV; and a per-host terms-of-service check before any
bulk run remains prudent.

\section{Minutes: how far back, and can we fetch them}

Four probes answer one question. \texttt{archive\_depth\_probe/archive\_depth.csv} (109 $\times$
10) is the post-merge authority on depth; \texttt{civicplus\_depth\_probe/civicplus\_depth.csv}
and \texttt{migration\_cliffs\_probe/migration\_cliffs.csv} feed corrections into it;
\texttt{minutes\_platform\_pilot/} is the only probe that touched an actual document.

\subsection{Depth distribution and the panel window}

\textbf{Pass 1} read each of the 109 archives as a normal browser client --- no downloads,
time-boxed, \emph{unknown over guess} --- and left 55 of 109 unknown, mostly because JS-shell
portals do not put depth in static HTML. \textbf{Pass 2} (\texttt{api\_probe.py}, 22 live jobs)
queried each portal's own API for metadata only: Legistar \texttt{events} filtered on
\texttt{EventMinutesFile ne null}, CivicClerk \texttt{/v1/Events} scanning
\texttt{publishedFiles} for \texttt{type="Minutes"}, PrimeGov \texttt{ListArchivedMeetings}, and
the Granicus minutes-mode RSS feed (capped at roughly 80--100 items). \textbf{The typing rule} is
what makes the numbers mean anything: a year counts only if the API returned an actual
Minutes-type file --- agendas, packets, video and year-dropdown ranges do not count.

\begin{center}
\begin{tabular}{lrr@{\qquad}lrr}
\toprule
\textbf{Earliest minutes year} & \textbf{Count} & \textbf{Share} &
\textbf{Common start} & \textbf{Localities} & \textbf{high/med} \\
\midrule
$\leq$ 2005      & 6  & 5.5\%  & $\leq$ 2014 & 29 & 22 \\
2006--2010       & 11 & 10.1\% & $\leq$ 2016 & 36 & 29 \\
2011--2015       & 15 & 13.8\% & $\leq$ 2018 & 45 & 37 \\
2016 and later   & 36 & 33.0\% &             &    &    \\
\texttt{unknown} & 41 & 37.6\% &             &    &    \\
\midrule
\textbf{Known depth} & \textbf{68} & \textbf{62.4\%} & & & \\
\bottomrule
\end{tabular}
\end{center}

A 2016-start panel therefore has enough co-submarket breadth for a strategic-interaction reaction
function around the ramp, and 36 is a \textbf{floor}: the 41 \texttt{unknown} rows are
overwhelmingly a method limit (JS or API-gated portals), not proven-shallow archives, so resolving
them can only add localities. Deep history is the opposite story --- only 6 localities reach 2005
or earlier and only 17 reach 2010 or earlier --- so long-run ``ratchet'' dynamics are not
supported by the online record. Supporting distributions: confidence \texttt{high} 32,
\texttt{med} 23, \texttt{low} 54; continuity \texttt{continuous} 48, \texttt{gaps} 3 (Cloverdale
2012 missing, Fairfax 2025 missing, Daly City sparse), \texttt{unknown} 58; type basis
\texttt{typed\_minutes} 42, \texttt{type\_ambiguous} 30, \texttt{agendas\_only} 2,
\texttt{unknown} 35.

Of the 55 rows in \texttt{archive\_depth\_api.csv} (55 $\times$ 11),
\texttt{minutes\_type\_confirmed} is \texttt{yes} for 20, \texttt{no\_api} for 34,
\texttt{absent} for 1. Nineteen resolved to a specific verified year: Alameda 2005; El Cerrito
and San Carlos 2008; Albany and Daly City 2009; Oakland and Burlingame 2014; Hayward, Emeryville
and Sonoma County 2015; Santa Rosa 2016; Hercules, Napa and San Mateo County 2017; Pinole and San
Mateo 2019; Napa County 2021; Contra Costa County and Pleasanton 2023. The twentieth ``yes'' is
the city of Sonoma, whose CivicWeb portal confirms a Minutes bucket labelled ``prior to 2017''
without exposing the year. The one \texttt{absent} is \textbf{Cupertino}: Legistar holds events
with status Final back to 2005 but \texttt{EventMinutesFile} is null throughout, so its minutes
live elsewhere. Platform mix of the 55: Legistar/Granicus 20, CivicPlus 14, CivicClerk 5, custom
CMS 4, PrimeGov 3, CivicWeb 3, Municode Meetings 2, and one each of eScribe, OnBase, NovusAgenda
and other.

\subsection{The agenda trap}

Year-dropdown floors systematically overstate depth, because dropdowns span \emph{agendas} or
events, which predate posted minutes. Hard checks against Legistar \texttt{EventMinutesFile}:

\begin{center}
\begin{tabular}{lccl}
\toprule
\textbf{Locality} & \textbf{Dropdown floor} & \textbf{Actual minutes} & \textbf{First minutes record} \\
\midrule
Oakland    & 2000 & \textbf{2014} & 2014-09-18 \\
Santa Rosa & 1999 & \textbf{2016} & 2016-02-02 \\
Hayward    & 2013 & \textbf{2015} & --- \\
Alameda    & 2005 & 2005 (confirmed) & 2005-01-12 \\
Emeryville & 2015 & 2015 (confirmed) & --- \\
\bottomrule
\end{tabular}
\end{center}

Oakland alone was a 14-year overstatement. Any remaining non-API dropdown floor among the 68
known rows is an \emph{upper bound} until API-confirmed.

\subsection{CivicPlus: the method works, the coverage does not}

CivicPlus AgendaCenter exposes no depth API, leaving 14 localities as the largest unresolved
block --- but it does have a server-rendered Search endpoint answering a date-range query in plain
HTML (\texttt{/AgendaCenter/Search/?term=\&CIDs=\{cid\}\&startDate=\dots\&endDate=\dots}) and
returning links of the form \texttt{ViewFile/Minutes/\_YYYYMMDD-N}. Because the document type is
\emph{in the path}, the read is definitively minutes-typed --- the property dropdown reads lacked,
and the reason the agenda trap cannot bite here. Regex the planning body's numeric \texttt{CID}
out of \texttt{/AgendaCenter}, query Search, take the minimum year in \texttt{ViewFile/Minutes}
links, collecting agenda years separately so an agenda/minutes gap is flagged rather than mistaken
for depth. The first pass used one wide 2004--2018 window; the second walked single probe years
(2008, 2010, 2012, 2014, 2016), avoiding truncation and giving a crude continuity check.

\textbf{Result: 2 of 14 resolved.} \textbf{Campbell 2006} is the win --- minutes-typed files
continuously 2006--2016 under \texttt{CID=6}, with 2004 empty: a clean verified floor, and one of
only six Bay Area localities reaching 2006 or earlier. \textbf{Los Altos Hills 2022} is a floor,
not a depth --- static \texttt{ViewFile/Minutes} links for 2022--2026 with the deeper archive
still gated, recorded at \texttt{med} confidence flagged
\texttt{static\_html\_floor(deeper\_may\_exist)}. The other twelve --- Antioch, Atherton, Cotati,
Half Moon Bay, Millbrae, Monte Sereno, Oakley, Pleasant Hill, San Anselmo, San Pablo, Windsor,
Yountville --- run the newer \textbf{CivicEngage} generation, which renders the committee list
client-side: the \texttt{CID} never appears in static HTML and \texttt{PreviousVersions} returns
404, so the query cannot be formed. That is a \emph{rendering} limit, not an access one (HTTP 200
throughout, hosts confirmed unblocked by the Akamai sweep), so all twelve are honest
\texttt{unknown} flagged \texttt{civicengage\_js\_gated(CID\_not\_static)} --- explicitly
\textbf{not} \texttt{minutes\_absent}. One browser-rendered session reading each \texttt{CID} out
of the rendered DOM finishes the job; it is cheap and not a blocker, since the twelve are mostly
small cities outside the estimation panel.

\subsection{Migration cliffs: a change of address, not a loss of data}

Three jurisdictions showed minutes apparently starting around \emph{2024} because they migrated
their portal that year --- and they are among the largest regulators in the set (San Jose is the
region's biggest city; the two counties regulate all unincorporated territory), so a 2024 start
would have excluded them from the ramp panel. Each prior platform was identified from redirects,
clerk pages and URL conventions, tested for whether the legacy host still serves, and read at
listing level using minutes-only feeds (Granicus \texttt{ViewPublisherRSS.php?mode=minutes}) or
year-filtered calendars (IQM2). \textbf{All three cliffs were domain or view migrations. 3 of 3
gain deep history.}

\textbf{Marin County --- 2005, high confidence}: Board of Supervisors agendas, minutes and video
have been online since 2005 on the legacy Granicus view (\texttt{view\_id=33}), and a separate
older archive at \texttt{marincounty.org/depts/bs/oldmeetingarchive} carries 1995--2005 agendas
and minutes without video. \textbf{Santa Clara County --- 2008, medium}: the legacy IQM2 instance
at \texttt{sccgov.iqm2.com/Citizens/} is still up and returns Board Minutes links back to 2008,
with the CSV explicit that the type is strongly indicated but \emph{not} document-verified and
that the archive likely extends earlier. \textbf{San Jose --- 2005 at the listing, documents
gated}: the legacy Granicus minutes RSS (\texttt{view\_id=51}) lists \textbf{80 minutes items
spanning 2005--2018}, but the MinutesViewer link 302-redirects to the \texttt{sanjoseca.gov}
DocumentCenter, which is Akamai-blocked; recorded as
\texttt{prior\_archive\_reachable = listing\_only} with an explicit instruction not to attempt to
defeat the block.

All three merged into \texttt{archive\_depth.csv}; together with Campbell 2006 and Los Altos
Hills 2022 they took the confirmed $\leq$2016 panel \textbf{from 32 to 36 localities}. The lesson
generalises --- a portal migration is a \emph{visible-archive cliff}, and any locality whose
depth begins at a suspiciously round recent year should be checked against its prior platform ---
but only these three were investigated, so other cliffs may sit undetected inside the 41
\texttt{unknown} rows. Both surviving legacy archives are decommissioning candidates, so
retrieving Marin and Santa Clara County history is time-sensitive.

\subsection{The retrieval and extraction pilot: the one probe that touched a document}

The earlier 14-city pilot (Daly City, South San Francisco, San Mateo, Redwood City, San Bruno,
Burlingame; San Jose, Palo Alto, Mountain View, Sunnyvale; Oakland, Berkeley, Fremont, Richmond)
is superseded in \emph{coverage} by the census but is the only source for the sprint's most
consequential finding.

\paragraph{No single-vendor scraper covers the set.} The 14 cities span six stacks:
Legistar/Granicus $\times 6$, CivicPlus $\times 4$, CivicClerk $\times 1$, PrimeGov $\times 1$,
Hyland OnBase $\times 1$, eScribe $\times 1$, plus Berkeley on a custom CMS with Laserfiche.
Berkeley and the City of San Mateo could not be cleanly classified, and two platforms present ---
eScribe and Hyland OnBase --- are unsupported by \texttt{civic-scraper} in any form.

\paragraph{One of three pilot cities yielded minutes.} \emph{Daly City} succeeded, but not
through \texttt{civic-scraper}, which targets the legacy CivicClerk ASPX portal while Daly City
runs the modern \texttt{*.api.civicclerk.com} React backend. The modern OData API does work: an
\texttt{Events} query filtered on \texttt{categoryName eq 'Planning Commission'} returns events
whose \texttt{publishedFiles[]} carries typed Agenda, Minutes and Agenda Packet entries,
downloadable via \texttt{GetMeetingFileStream} --- the pattern later reused for the roughly ten
CivicClerk localities region-wide. \emph{San Jose} failed on every route: its Legistar Web API
returned complete metadata (14 Planning Commission meetings for the first half of 2024, matching
the second-and-fourth-Wednesday schedule) but every one has \texttt{EventMinutesFile = False} and
status ``Draft'', and the Granicus route dead-ends at the Akamai-blocked DocumentCenter.
\emph{Fremont} failed on access: \texttt{civic-scraper} returned \textbf{zero assets}, and
\texttt{/AgendaCenter}, \texttt{/government/agenda-center} and the RSS feed all return HTTP 403.

\paragraph{A silent-gap risk: publishing minutes is not the same as holding meetings.} Of the
\textbf{12} Daly City Planning Commission meetings in 2024, the API exposed Agenda $\times 11$,
Agenda Packet $\times 5$, Minutes $\times 1$ --- \emph{one} minutes document for the year, so an
unattended backfill would report success and return essentially nothing. Alongside it, minutes
are approved one meeting later and attached to the \emph{following} meeting's packet (the 5
December 2023 minutes arrived under the 6 February 2024 packet), so a backfill must key minutes
to meeting $N-1$. The same pattern recurs in the census, where Campbell and Healdsburg show
agendas but no posted minutes.

\paragraph{The central finding: the SF extractor does not transfer.} Run unmodified on the real
Daly City minutes, \texttt{autoextract.extract()} produced, out of roughly 35 schema fields,
\textbf{two} correct ones --- \texttt{ceqa\_determination} (\texttt{eir}; an FEIR was before the
commission) and \texttt{item} (1). The rest were empty or confidently wrong:

\begin{center}
\small
\begin{tabular}{p{0.17\textwidth} p{0.15\textwidth} p{0.24\textwidth} p{0.32\textwidth}}
\toprule
\textbf{Field} & \textbf{Extractor} & \textbf{Truth} & \textbf{Failure mode} \\
\midrule
\texttt{case\_number} & \texttt{21-14998} & \texttt{GPA-04-21-14998} &
SF regex grabbed the tail of the case number \\
\texttt{units\_proposed} & \texttt{235} & \texttt{1,235} &
the thousands comma split the match \\
\texttt{parking\_spaces} & \texttt{5} & \texttt{34.5} &
the decimal point split the match \\
\texttt{project\_address} & \texttt{333 90Th Street} & city-hall boilerplate &
grabbed the page footer \\
\texttt{request\_type} & empty & GPA $+$ PD $+$ subdivision $+$ design review &
SF uses case-number \emph{suffixes}; Daly City uses \emph{prefixes} \\
\texttt{action} & empty & ``voted 4-0 \dots\ recommend approval'' &
no SF-style \texttt{ACTION:} label in the document \\
\texttt{assessor\_block} & empty & \texttt{APN 091-211-230} &
Daly City writes APN, not ``Assessor's Block'' \\
vote, speakers, stance & empty / 0 & 4-0 vote; no stance markers &
SF's post-2015 stance markers do not exist here \\
\bottomrule
\end{tabular}
\end{center}

The dangerous rows are the middle three: an empty field announces itself; a populated address,
unit count and parking count that are all silently wrong do not. The operational conclusion,
carried forward as a \emph{gating quality step}: the SF heuristics must never run unattended on
another jurisdiction, and any multi-jurisdiction build must first validate an extractor on real
Granicus or CivicClerk minutes \emph{with human review}.

\paragraph{Vocabulary does not transfer either.} The proposed mapping layer
(\texttt{jurisdiction\_mappings.py}, additive, not wired into the pipeline) covers Daly City
only: of its 13 \texttt{request\_type} entries, five are clean synonyms (general plan amendment,
GPA, rezoning, conditional use, variance) and eight are flagged judgment calls; of seven
\texttt{action} entries, five are clean and two flagged. Three flagged cases matter
substantively. \textbf{``DR'' means different things in different cities}: in SF it is
\emph{Discretionary Review} --- review of a project that is otherwise by right, precisely the
behaviour the model is about --- while in Daly City it is \emph{Design Review}, so mapping on the
letters would be exactly wrong. \textbf{Four common suburban request types have no canonical
enum} --- Planned Development, Major Subdivision, Precise Plan, Development Agreement --- and all
collapse to \texttt{other}; they recur across California suburbs, so it is a schema question. And
\textbf{``recommend approval'' is not approval}: on legislative items the Daly City commission
only recommends to the City Council, while the SF \texttt{action} enums assume a final
disposition. The pilot document is exactly that case --- a 4-0 vote to \emph{recommend} approval
of a general plan amendment, rezoning, subdivision, design review and EIR certification for a
\textbf{1,235-unit} project.

\paragraph{Tooling limits, recorded because they will recur.} \texttt{civic-scraper} v1.1.0 has
four defects that bit here: its CLI auto-routes only CivicPlus and DigitalTowPath;
\texttt{LegistarSite} crashes in \texttt{pytz.timezone(None)} unless a timezone is passed; its
Legistar events call fetches the current year only (\texttt{events(since=2024)} returned 180
events all dated 2026), so backfill needs the Legistar Web API with date filters; and its
CivicClerk and CivicPlus scrapers target older portal versions than several Bay Area cities run.

\section{Zoning: the current envelope and the pre-period envelope}

\subsection{Spatial form: a script, not an army of RAs}

\texttt{zoning\_map\_form\_probe/zoning\_map\_form.csv} (109 $\times$ 8) classifies the
\emph{form} in which each locality publishes its current spatial zoning --- deliberately not what
the zoning says --- looking in order for a downloadable GIS layer, a viewer with no obvious
download, a static PDF map, or nothing.

\begin{center}
\begin{tabular}{lrrl}
\toprule
\textbf{\texttt{best\_form}} & \textbf{Count} & \textbf{Share} & \textbf{Meaning} \\
\midrule
\texttt{gis\_layer}        & 67 & 61.5\% & downloadable shapefile / GeoJSON / feature service \\
\texttt{gis\_viewer\_only} & 28 & 25.7\% & interactive viewer, no obvious download \\
\texttt{pdf\_map}          & 13 & 11.9\% & static zoning-map PDF only \\
\texttt{none\_found}       & 1  & 0.9\%  & nothing located (Pacifica) \\
\bottomrule
\end{tabular}
\end{center}

\texttt{download\_apparent}: yes 71, no 37, unknown 1. Confidence: high 60, med 46, low 3
(Hillsborough, San Bruno, Pacifica). The 13 PDF-only localities --- Atherton, Clayton, Colma,
Daly City, Healdsburg, Hercules, Millbrae, Monte Sereno, Pinole, Portola Valley, San Pablo,
Saratoga, Woodside --- are all small, none a large regulator. Within \texttt{gis\_layer},
\textbf{55 are native} city, county or consortium sources and \textbf{12 are covered only by the
statewide fallback}: twelve Alameda County cities (Alameda, Albany, Dublin, Emeryville, Fremont,
Hayward, Livermore, Newark, Piedmont, Pleasanton, San Leandro, Union City) whose own portals were
never separately checked --- verifiable from the file, since exactly 12 \texttt{gis\_layer} rows
carry a \texttt{gis.data.ca.gov} URL. Three regional or county sources close 26 localities on
their own: MarinMap (all 12 Marin localities, 11 rows via \texttt{gis.marinpublic.com} plus the
county's own), the Solano Regional GIS Consortium (all 8 Solano localities on one ArcGIS
organisation), and Napa County GIS (all 6 Napa localities, 5 on \texttt{gis.napacounty.gov} plus
the city of Napa's own). Contra Costa, San Mateo, Santa Clara and Sonoma county GIS cover
unincorporated areas only; their cities largely self-publish. Above it all sits the CA Statewide
Zoning layer (Gov-OPR, \texttt{gis.data.ca.gov}), one downloadable layer plus REST, reported as
covering \textbf{535 of 539} California jurisdictions.

Spatial matching is therefore a \textbf{scriptable problem}: native layers where they exist, the
statewide layer as gap-filler for the 28 viewer-only, 13 PDF-only and 1 none-found localities, so
no locality is a dead end for automation --- decisively better than the minutes side, where
roughly 38\% still have unverified depth. \textbf{The one substantive risk} is that
the statewide layer is a \textbf{2022--23 aggregated snapshot} of varying per-city fidelity,
underpinning automation for about \textbf{54 of 109} rows once the 12 unchecked Alameda cities
are counted; since the by-right envelope is the load-bearing right-hand-side variable, it should
be spot-checked against each city's own viewer or PDF --- accept it as the universal first cut,
or budget native extraction for the roughly 42 localities that natively offer only a viewer or a
PDF. Open, cheapest first: confirm feature-service endpoints behind the 28 viewer-only sites (the
true \texttt{gis\_layer} count is likely above 67; the exceptions are El Cerrito, Lafayette and
Martinez on the non-ArcGIS ``CommunityView'' product); check the 12 Alameda cities natively;
resolve the three low-confidence identifications. Caught in passing: the brief's claim that ``20
rows already mention GIS'' in the census was a substring artefact --- the ``gis'' inside
Le\emph{gis}tar --- so all 109 were probed fresh. Noted but not acted on: Stanford EarthWorks
holds a \textbf{historical 2014 version of the San Mateo County planning-zones layer}, a pointer
for the deferred historical-envelope workstream.

\subsection{What the code says: NZLUD, code hosts, and the extraction pilot}

\texttt{zoning\_envelope\_project/} worked the same 14-city set as the minutes pilot, asking
whether a machine-readable by-right dataset exists, whether the envelope can be read out of
municipal codes at acceptable cost, and whether the preemption shock can be dated: mostly no, yes
but not mechanically, and yes cleanly.

\paragraph{NZLUD covers 4 of 14, and only pre-preemption.} The National Zoning and Land Use
Database (Mleczko and Desmond, 2023) is an NLP-derived reading of municipal codes collected
2019--2022, replicating the 2006 Wharton WRLURI frame plus 210 supplemental municipalities; the
full release in \texttt{\_source\_data/nzlud\_muni.csv} is 2,639 municipalities $\times$ 77
columns, and its key field \texttt{mf\_per} is the proportion of residential districts (including
mixed-use) permitting multifamily housing \emph{by right} --- close to the margin the model
needs. Matching on Census place FIPS yields exactly four rows. The ten absent cities include most
of the Peninsula and every East Bay city except Fremont --- Daly City, South San Francisco, San
Mateo, San Bruno, Palo Alto, Mountain View, Sunnyvale, Oakland, Berkeley, Richmond --- a direct
consequence of the WRLURI-2006 frame, which no amount of waiting will fix.

\begin{center}
\begin{tabular}{lrrrrl}
\toprule
\textbf{City} & \texttt{mf\_per} & \texttt{zri} & \texttt{height\_ft\_median} &
\texttt{parking\_median} & \textbf{Snapshot} \\
\midrule
Burlingame   & 0.667 & 2.925 & 50 & 2 & 2021-02-19 \\
Redwood City & 0.600 & 2.533 & 50 & 1 & 2019-10-24 \\
Fremont      & 0.467 & 2.537 & 30 & 1 & 2019-10-15 \\
San Jose     & 0.389 & 3.318 & 45 & 1 & 2021-12-19 \\
\bottomrule
\end{tabular}
\end{center}

Every snapshot predates the 2023--2026 sixth-cycle Housing Element updates and predates or
coincides with the SB-9, SB-35 and SB-423 rezonings, so NZLUD is best read as a
\emph{pre-preemption} envelope: a limitation for measuring today's rules, an opportunity for a
two-period design. The one available spot-check is close but not identical --- Fremont's Title 18
codes multifamily as \texttt{P} in R-3 and R-G and prohibited in R-1 and R-2, a naive
base-residential by-right share of $2/4 = 0.50$ against NZLUD's $0.467$, the gap being a
denominator disagreement (NZLUD counts mixed-use districts as residential) rather than an obvious
error in either. San Jose could not be checked at all because Municode blocked the read; its
$0.389$ is unverified, and its Title 20 uses an \texttt{S} $=$ Special Use Permit symbol easy to
misread as by-right.

\paragraph{Half the code hosts block bots.} All 14 codes were located: Municode $\times 6$,
eCode360 $\times 3$, public.law $\times 2$, QCode $\times 2$, American Legal $\times 1$,
CodePublishing $\times 1$ (South San Francisco is ambiguous between QCode and eCode360 and is
counted in both). Only \textbf{CodePublishing and public.law} returned real code text to a plain
GET: Municode returns HTTP 200 with a roughly 6\,KB Angular shell and zero code text (the
ordinance loads from a separate JSON API), while American Legal, eCode360 and Berkeley's host
return 403 outright. This matches the NZLUD authors' own note that they could not use web
scraping and downloaded codes manually.

\paragraph{Two structural traps, both fatal to a naive parser.} Fremont's code is a consolidated
use matrix defining its own symbols (\texttt{P} permitted, \texttt{C} conditional use permit,
\texttt{Z} zoning-administrator permit, \texttt{--} prohibited) and parsed cleanly. San Mateo
splits permitted uses (\S27.22.010) from special uses (\S27.22.020) into \emph{separate sections
per district}, which a generic table parser would miss entirely; worse, \S27.22.055 says building
height ``shall not exceed the standards set forth on the Building Height Plan of the General
Plan'' --- \textbf{the height limit is not in the code text at all}. Height cannot be
text-scraped for such cities, which also makes NZLUD's height fields unreliable wherever a code
defers to a map. Both pilot cities carry unresolved overlay districts (Fremont: TOD, hillside,
housing-element-sites; San Mateo: two-unit and mixed-use residential).

\paragraph{The WRLURI benchmark disagrees with NZLUD.} \texttt{wrluri\_crosscheck.csv} (14 rows
$\times$ 8 columns) shows NZLUD covering 4 of 14, WRLURI 2006 covering the \emph{same} 4 of 14
(confirming the inherited frame), WRLURI 2018 covering 1 of 14 usable (Sunnyvale, 0.63), and all
three overlapping on \textbf{zero} cities. Over the four-city overlap the indices are
rank-inverted: Spearman between WRLURI 2006 and NZLUD \texttt{zri} is $-1.00$ where the expected
sign is positive, and between WRLURI 2006 and \texttt{mf\_per} is $+0.40$ where the expected sign
is negative, with San Jose driving it --- least restrictive of the four by the 2006 survey, most
restrictive by NZLUD's code index. This is neither validation nor refutation: $n=4$ makes a
$\pm 1.00$ correlation nearly meaningless, the vintages are 13--15 years apart across
California's entire preemption shift, and the constructs differ by design. The takeaway is a
caution --- per-cell measures such as \texttt{mf\_per} and height are more defensible to carry
forward than either composite index.

\subsection{The pre-period envelope: 15 of 25 recoverable}

The difference-in-differences needs a \emph{before} picture of the envelope, but codes and zoning
maps are published as \emph{current} documents, silently overwritten. The Wayback probe asks only
whether a circa-2016 envelope can be \emph{located and dated} cheaply across the 25-locality
sample. In-window means a Wayback CDX capture timestamped \textbf{2014--2018}, and the capture
timestamp \emph{is} the vintage evidence --- the probe never infers a date from content. Two forms
count: an archived municipal-code page (text form, whose use table is the envelope) and an
archived zoning-map PDF (spatial form). Deep code URLs were broadened to the publisher's
code-library root so prefix matching catches any in-window capture, and basenames were verified by
hand because shared publisher hosts produce cross-jurisdiction matches. The standing rule is
\texttt{none\_found} over a guess.

\textbf{A usable in-window pre-period source exists for 15 of 25 localities (60\%):} 11 archived
code-page captures, 3 archived zoning-map PDFs, and San Francisco (in hand via the project's own
pipeline plus DataSF historical zoning layers). The remaining 10 are 9 \texttt{none\_found} plus
Burlingame, whose only source is an NZLUD 2019--21 proxy that is post-SB-35 and therefore not a
valid pre-period at all. \texttt{confidence} reads high for exactly 15 rows and low for 10, and
the CSV contains \textbf{zero live false-positive source URLs}. Rejected in the first pass: an
Oro Valley, Arizona PDF on the shared \texttt{codepublishing.com} host that had polluted both
Solano County and Walnut Creek; Palo Alto's \texttt{Symposium.pdf}; and an SF SoMa area guide.
Two more were killed in a second pass --- \textbf{San Ramon}, whose 2015 \texttt{rezoning.pdf} is
a seven-page ``Rezoning to P-1 Submittal Requirements'' \emph{application}, not the citywide map
(a lead survives: a \texttt{zoning\_ordinance\_map} page captured 2018, but it 302-redirects to
EncodePlus and the underlying map's vintage is unverified); and \textbf{Daly City}, whose 2014
``Zoning Ordinance.pdf'' sits under a General-Plan-Update (\texttt{/gpu/}) path and could not be
confirmed adopted (Wayback returned 404/503). Both were reset to \texttt{none\_found}. Solano
County (2015) and Walnut Creek (2014) survived as legitimate CodePublishing code-page captures;
Palo Alto had already been reset.

The three surviving spatial anchors are the cleanest vintages --- Sunnyvale 2014, Tiburon 2015 (a
map effective 2006-03-31 still posted in-window, so the circa-2016 envelope was stable), San
Leandro 2016 --- while the code-page captures are 2014 (San Carlos, Walnut Creek), 2015 (Solano
County), 2017 (Alameda, Mountain View, Oakland, Sonoma County) and 2018 (Clayton, El Cerrito,
Hayward, Milpitas). That 2017--2018 cluster is comfortably before SB-9, SB-423 and the builder's
remedy but \emph{after} SB-35 (2017), so for an SB-35-specific margin those are ``early-ramp''
rather than true pre-period and the 2014--2016 PDFs are the cleaner anchors. \textbf{The gap is
patterned, not random}: the 9 \texttt{none\_found} localities cluster by code publisher ---
eCode360 (Albany, Emeryville, Santa Rosa), American Legal \emph{legacy} platform (Palo Alto,
Fairfax, whose 2016 code sits on \texttt{library.amlegal.com} while the probe queried
\texttt{codelibrary.amlegal.com}), public.law (Petaluma), a city-site PDF case (Piedmont), plus
the two killed above --- so the fix is a per-publisher amendment/superseded-version dig, a bounded
build task. But the whole probe is downstream of the unmade envelope-definition decision: located
sources are worth little until the extraction target is defined.

\section{The treatment variable: builder's-remedy exposure}

\texttt{hcd\_preemption\_panel/hcd\_preemption\_panel.csv} (25 $\times$ 10) is the ``preemption
ramp'' variable that dates the shock and supplies the instrument for the reaction-function
estimation. The measure is exposure to California's \textbf{builder's remedy}: while a locality
lacks a Housing Element found in substantial compliance by HCD it cannot deny qualifying
affordable projects on zoning or general-plan grounds, so exposure runs from the statutory
adoption deadline to the HCD compliance finding --- rule-driven timing plausibly orthogonal to a
locality's contemporaneous demand shock.

The authoritative input is HCD's \emph{Housing Element Review and Compliance Report}
(\url{https://data.ca.gov/dataset/housing-element-compliance-report}), pulled unmodified to
\texttt{\_source\_data/hcd.csv} --- 539 rows $\times$ 10 columns, one per California jurisdiction,
last updated 2026-06-05 --- filtered to \texttt{Record Type} $=$ \emph{Adopted}; prohousing status
comes from HCD's \emph{Prohousing Designated Jurisdictions} list as of 2026-03-26. The sample is
the 25 localities with a verified minutes start year of 2016 or earlier at high/medium confidence.
Two date conventions are deliberately not collapsed: \texttt{he\_adoption\_date} (HCD's \emph{Date
Received}, a proxy for self-adoption) and \texttt{hcd\_certification\_date} (HCD's \emph{Reviewed
Date}, the formal finding). With the ABAG sixth-cycle deadline $T_0 = \text{2023-01-31}$,
\[
  \text{exposure}_j \;=\; \max\!\big\{0,\; \text{months}\big(T_0 \rightarrow
  \text{HCD compliance finding}_j\big)\big\},
\]
with $\max\{0, \text{months}(T_0 \rightarrow \text{HE received}_j)\}$ recorded alongside;
localities certified before the deadline are floored at zero, not given negative exposure.

\paragraph{The variable is clean and the variation is strong.} All 25 matched an HCD
adopted-element record, all 25 carry \texttt{confidence = high} and sourced dates, and there are
no \texttt{no\_hcd\_match} rows. Exposure ranges from \textbf{0.0 to 38.4 months} (mean 9.9,
median 8.6): three localities have zero exposure, 22 positive, nine exceed 12 months and two
exceed 24 months. 13 of 25 are prohousing-designated; five have sourced designation dates
(Mountain View and Petaluma January 2024, Walnut Creek 2024-08-01, Santa Rosa and Sonoma County
2025) and the other eight carry \texttt{prohousing\_date\_to\_verify(HCD tracker XLS)}.

\begingroup\small
\begin{longtable}{@{}lccrl@{}}
\caption{Builder's-remedy exposure, 25-locality estimation sample (deadline 2023-01-31
$\rightarrow$ HCD compliance finding). ``designated'' $=$ prohousing-designated, exact date to
verify; ``---'' $=$ not designated.}\\
\toprule
Locality & HE received & HCD finding & Exposure (mo.) & Prohousing \\
\midrule
\endfirsthead
\toprule
Locality & HE received & HCD finding & Exposure (mo.) & Prohousing \\
\midrule
\endhead
\bottomrule
\endfoot
Alameda        & 2022-11-16 & 2022-12-20 & 0.0  & designated \\
San Francisco  & 2023-01-31 & 2023-02-01 & 0.0  & designated \\
San Leandro    & 2022-12-09 & 2023-02-02 & 0.0  & designated \\
Emeryville     & 2023-01-20 & 2023-02-03 & 0.1  & designated \\
Oakland        & 2023-02-13 & 2023-02-17 & 0.5  & designated \\
San Ramon      & 2023-02-06 & 2023-02-27 & 0.9  & --- \\
Milpitas       & 2023-05-16 & 2023-05-17 & 3.5  & --- \\
Petaluma       & 2023-03-22 & 2023-05-18 & 3.6  & 2024-01 \\
Mountain View  & 2023-04-26 & 2023-05-26 & 3.8  & 2024-01 \\
Hayward        & 2023-06-07 & 2023-07-27 & 5.9  & designated \\
El Cerrito     & 2023-08-17 & 2023-08-22 & 6.7  & designated \\
Albany         & 2023-08-16 & 2023-09-08 & 7.2  & --- \\
Tiburon        & 2023-09-22 & 2023-10-18 & 8.6  & --- \\
Walnut Creek   & 2023-08-25 & 2023-10-24 & 8.8  & 2024-08-01 \\
Sonoma County  & 2023-08-30 & 2023-10-26 & 8.8  & 2025 \\
Piedmont       & 2023-09-11 & 2023-11-09 & 9.3  & --- \\
Sunnyvale      & 2024-02-29 & 2024-03-06 & 13.2 & designated \\
Burlingame     & 2024-03-15 & 2024-03-20 & 13.6 & --- \\
Fairfax        & 2024-03-14 & 2024-04-08 & 14.2 & --- \\
Solano County  & 2024-02-09 & 2024-04-09 & 14.3 & --- \\
San Carlos     & 2024-04-19 & 2024-04-25 & 14.8 & --- \\
Palo Alto      & 2024-07-26 & 2024-08-20 & 18.6 & --- \\
Daly City      & 2024-11-13 & 2024-12-03 & 22.1 & --- \\
Santa Rosa     & 2025-07-28 & 2025-09-12 & 31.4 & 2025 \\
Clayton        & 2026-02-12 & 2026-04-13 & 38.4 & --- \\
\end{longtable}
\endgroup

\paragraph{The earlier 14-city panel is not a subset of this one.}
\texttt{zoning\_envelope\_project/hcd\_preemption\_exposure\_panel.csv} holds 14 rows, all
planning period \texttt{6S}, record type \texttt{Adopted}, current compliance \texttt{In},
anchored at \textbf{2023-02-01} rather than 2023-01-31, with exposure spanning \textbf{0.5 to
40.0 months}: Oakland 0.5 (found compliant 2023-02-17), Berkeley 0.9 (2023-02-28), Fremont 1.6
(2023-03-22), Redwood City 1.8 (2023-03-27), Mountain View 3.8 (2023-05-26), Richmond 8.0
(2023-10-02), South San Francisco 9.6 (2023-11-20), San Jose 11.9 (2024-01-29), Sunnyvale 13.1
(2024-03-06), Burlingame 13.6 (2024-03-20), Palo Alto 18.6 (2024-08-20), San Bruno 20.7
(2024-10-21), Daly City 22.1 (2024-12-03), San Mateo 40.0 (2026-06-01). Eight localities ---
Berkeley, Fremont, Redwood City, Richmond, San Bruno, San Jose, San Mateo and South San Francisco
--- appear \emph{only} in this 14-row file, the two files use different column schemas, and they
anchor exposure differently, so neither should be deleted and the anchor convention must be
unified before they are used together. San Mateo is the striking case: roughly 3.3 years of
continuous exposure, certifying only in mid-2026 while its Peninsula neighbours cleared in months
--- an attractive treatment case, provided the minutes side can retrieve its hearings, which the
minutes pilot could not even confidently classify.

\paragraph{Caveats.} The \textbf{snapshot caveat} is the main analytical limitation: the HCD
dataset is a \emph{current-status} snapshot and all 25 read as ``In'' compliance today, so any
decertification-then-recertification sequence during the ramp is invisible and exposure would be
understated; \texttt{status\_sequence} flags this on every row, and verifying contested cases is
the one substantive \texttt{to\_verify}. At the full ABAG level \texttt{hcd.csv} shows 108 In and
one Out (Half Moon Bay, reviewed 2026-02-02). Also open: whether exposure ends at self-adoption
or at HCD's finding (the gap is material --- Solano County 2024-02-09 vs.\ 2024-04-09); the eight
unsourced prohousing dates; the 25-versus-36 sample width; and a possible \emph{second} preemption
series from HCD's SB-35/SB-423 Statutory Determinations, RHNA-progress-based and not built here.

\section{Data inventory}

Bulk data lives on Dropbox; everything here is a small derived table, a script, or a narrative
report. Status is \emph{current} (believe it), \emph{provenance} (kept for the record, do not
cite), \emph{superseded}, \emph{input} (unmodified third-party), or \emph{broken} (will not run
as written; \S7). Rows are grouped by the probe directory that produced them.

\begingroup\small
\begin{longtable}{@{}p{0.30\textwidth} p{0.41\textwidth} p{0.22\textwidth}@{}}
\toprule
\textbf{File} & \textbf{What it is} & \textbf{Status} \\
\midrule
\endfirsthead
\toprule
\textbf{File} & \textbf{What it is} & \textbf{Status} \\
\midrule
\endhead
\bottomrule
\endfoot
\multicolumn{3}{@{}l}{\textbf{\texttt{bay\_area\_census/}} --- the frame} \\
\texttt{bay\_area\_locality\_census.csv} & 109 $\times$ 20; platform, URL, access, solution class,
tier, confidence, both sides & \textbf{current} (the spine) \\
\texttt{raw/*.json}, \texttt{ca\_place\_geoid.json} & 8 per-county surveys (108 localities);
Gazetteer name $\rightarrow$ GEOID & provenance / current \\
\texttt{consolidate.py} & Builds the CSV, appends SF, attaches GEOIDs & current \\
\texttt{bay\_area\_locality\_census\_report.md} & Narrative roll-up, 2026-06-08 & provenance;
differs from CSV in 4 places \\
\texttt{akamai\_civicplus\_probe\_finding.md} & The 28-host CivicPlus sweep table & current \\
\texttt{sample\_saratoga\_pc\_}\newline\texttt{minutes\_2026-04-08.pdf} & 19\,KB; the one document
retrieved end-to-end & evidence \\
\addlinespace
\multicolumn{3}{@{}l}{\textbf{\texttt{archive\_depth\_probe/}} --- the time dimension} \\
\texttt{archive\_depth.csv} & 109 $\times$ 10; post-merge authority on minutes depth &
\textbf{current} \\
\texttt{archive\_depth\_api.csv} & 55 $\times$ 11; API depth-read detail & current \\
\texttt{consolidate.py} & Builds the CSV from \texttt{raw/} $+$ census GEOIDs & current;
``verified'' filter is buggy \\
\texttt{api\_probe.py} & Live API client (Legistar, CivicClerk, PrimeGov, Granicus RSS) & broken
(\texttt{/tmp/} I/O) \\
\texttt{assemble\_api.py} & Merges 19 verified years in; results transcribed literally & current \\
\texttt{raw/*.json} & 8 per-county shallow-probe observations & provenance \\
\texttt{archive\_depth\_report.md} & Shallow-probe narrative; carries a banner & superseded \\
\texttt{archive\_depth\_api\_report.md} & API-pass narrative; \emph{no} banner & superseded (says
64/32, not 68/36) \\
\addlinespace
\multicolumn{3}{@{}l}{\textbf{\texttt{civicplus\_depth\_probe/}} --- AgendaCenter depth} \\
\texttt{civicplus\_depth.csv} & 14 $\times$ 8; 2 resolved, 12 honest \texttt{unknown} &
\textbf{current} \\
\texttt{final\_pass.py} & Single-year Search probes; supersedes \texttt{probe.py} & current \\
\texttt{probe.py} & First pass, wide 2004--2018 query & broken (hard-coded
\texttt{/Users/danpost/}) \\
\texttt{\_results.json} & Raw first-pass output & provenance; contradicts CSV twice \\
\addlinespace
\multicolumn{3}{@{}l}{\textbf{\texttt{migration\_cliffs\_probe/}} --- the 2024 portal moves} \\
\texttt{migration\_cliffs.csv} & 3 $\times$ 8; legacy platform, reachability, legacy depth &
\textbf{current} (no script) \\
\addlinespace
\multicolumn{3}{@{}l}{\textbf{\texttt{minutes\_platform\_pilot/}} --- the 14-city pilot} \\
\texttt{minutes\_platform\_pilot\_report.md} & Platform table, scrape results, transfer test, 13
review items & current on transfer; superseded on coverage \\
\texttt{pilot\_scrape.py} & Daly City CivicClerk OData retrieval; stdlib only & current \\
\texttt{jurisdiction\_mappings.py} & Proposed additive vocabulary layer; not wired in & current;
path constant is depth-sensitive \\
\texttt{sample\_dalycity\_pc\_}\newline\texttt{2023-12-05\_minutes.txt} & 4.4\,KB; the document the
transfer test rests on & evidence \\
\addlinespace
\multicolumn{3}{@{}l}{\textbf{\texttt{zoning\_map\_form\_probe/}} --- spatial form} \\
\texttt{zoning\_map\_form.csv} & 109 $\times$ 8; form, URL, source type, download flag &
\textbf{current} \\
\texttt{consolidate.py}, \texttt{raw/*.json} & Assembles the CSV; 8 per-county form probes &
current / provenance \\
\texttt{zoning\_map\_form\_report.md} & Narrative; no banner, numbers agree with the CSV & stale
in status only \\
\addlinespace
\multicolumn{3}{@{}l}{\textbf{\texttt{preperiod\_envelope\_probe/}} --- the 2016 envelope} \\
\texttt{preperiod\_envelope.csv} & 25 $\times$ 9; 15 usable, zero live false positives &
\textbf{current} \\
\texttt{wayback\_probe.py}, \texttt{\_wayback\_results.json} & Wayback CDX query 2014--2018; raw
output, 25 records & current / provenance \\
\texttt{finalize.py} & Applies hand corrections and writes the CSV & broken (regenerates 17/25) \\
\texttt{preperiod\_envelope\_report.md} & Original write-up; carries a correction note &
superseded on 17/25 \\
\addlinespace
\multicolumn{3}{@{}l}{\textbf{\texttt{hcd\_preemption\_panel/}} --- the treatment variable} \\
\texttt{hcd\_preemption\_panel.csv} & 25 $\times$ 10; exposure 0.0--38.4 months & \textbf{current,
load-bearing} \\
\texttt{build\_hcd\_panel.py} & Deterministic builder; prints the exposure summary & broken
(\texttt{/tmp/} sample list) \\
\addlinespace
\multicolumn{3}{@{}l}{\textbf{\texttt{zoning\_envelope\_project/}} --- the 14-city assessment} \\
\texttt{zoning\_envelope\_assessment.md} & NZLUD coverage, 14-city code table, extraction pilot,
panel, 15 review items & superseded on coverage; one claim contradicted \\
\texttt{nzlud\_14city\_subset.csv} & The 4 matching NZLUD rows $\times$ 19 columns & current \\
\texttt{hcd\_preemption\_exposure\_panel.csv} & 14 rows; exposure 0.5--40.0 mo., anchored
2023-02-01 & current; \emph{not} a subset of the 25-row panel \\
\texttt{pilot\_extract.py} & Fremont/San Mateo code read; documents the Municode block & current
(network-dependent) \\
\texttt{wrluri\_crosscheck.py} / \texttt{.csv} & 14 $\times$ 8; NZLUD vs.\ WRLURI 2006/2018 & CSV
current; script broken \\
\addlinespace
\multicolumn{3}{@{}l}{\textbf{\texttt{\_source\_data/}} --- unmodified third-party releases} \\
\texttt{nzlud\_muni.csv} & Full NZLUD: 2,639 municipalities $\times$ 77 columns & input \\
\texttt{hcd.csv} & Full HCD compliance report: 539 rows $\times$ 10 columns, updated 2026-06-05 &
input \\
\end{longtable}
\endgroup

\noindent Several scripts reach sideways with \texttt{HERE.parent / "sibling"}
(\texttt{archive\_depth\_probe/consolidate.py}, \texttt{archive\_depth\_probe/assemble\_api.py}
and \texttt{zoning\_map\_form\_probe/consolidate.py} $\rightarrow$ \texttt{bay\_area\_census/};
\texttt{hcd\_preemption\_panel/build\_hcd\_panel.py} $\rightarrow$ \texttt{\_source\_data/}), so
moving this group as a unit is safe and splitting it is not.
\texttt{minutes\_platform\_pilot/jurisdiction\_mappings.py} additionally counts \texttt{.parent}
up to the repo root and must be re-counted on any change of depth --- when that constant is wrong
the import of the canonical enums fails silently and the validation degrades to a no-op.

\section{Known contradictions and stale artifacts}

The most actionable section in the sprint: places where two files in the repository disagree, or
where a script would not reproduce the file it supposedly builds. None has been silently fixed.

\subsection{Artifacts that contradict each other}

\begin{enumerate}

\item \textbf{The archive-depth \texttt{.md} reports understate their own CSV.}
\texttt{archive\_depth\_api\_report.md} says known depth is \textbf{64 of 109} and $\leq$2016
breadth is \textbf{32}, whereas \texttt{archive\_depth.csv} gives \textbf{68} and \textbf{36}
after the CivicPlus and migration-cliff merges --- and it carries \emph{no} supersession banner.
\texttt{archive\_depth\_report.md} is the older shallow snapshot (55 unknown, 27 verified, Marin
County at 2024, gaps 2) and does carry one. Believe the CSV.

\item \textbf{\texttt{consolidate.py}'s ``verified'' filter wrongly rejects API-verified rows.} It
defines the verified subset by rejecting any row whose \texttt{flags} contain
\texttt{minutes\_type\_unverified} or \texttt{oldest\_year\_minutes\_unverified}. The API merge
added \texttt{api\_verified} to those rows but never cleared the old shallow-probe flags, so the
filter now rejects rows that were in fact document-verified --- Alameda 2005, Oakland 2014,
Hayward 2015 and Santa Rosa 2016 among them. Re-running it yields \textbf{38 of 109} ``verified''
rows, an artefact of stale flags rather than a real count. \emph{Use the \texttt{confidence}
column, not the flag filter.}

\item \textbf{\texttt{civicplus\_depth\_probe/\_results.json} says Campbell 2018; the CSV says
2006.} The JSON records ``21 Minutes files, years 2018--2018'' from the wide 2004--2018 query.
The CSV's 2006 is the later and correct value --- the wide query returned a truncated or
most-recent slice, which is exactly why the second pass switched to single-year probes, and the
year-by-year hits support 2006. The same JSON also records Los Altos Hills as \texttt{unknown}
with \texttt{cid\_not\_found} where the CSV records 2022 from the static-HTML floor. Both stale
values remain in the JSON.

\item \textbf{Millbrae's flag contradicts its own method field.} Its \texttt{method} says
``CID=8 found but Search returned no minutes 2008--2016'' --- the CID \emph{was} discoverable and
the window came back empty --- but its \texttt{flags} value is
\texttt{civicengage\_js\_gated(CID\_not\_static)}, which describes the other eleven. Millbrae is
a different failure mode and should be re-probed on its own terms, including whether
\texttt{CID=8} is actually the planning body.

\item \textbf{The census report disagrees with the census CSV in four places.} Real discrepancies,
not rounding; the CSV is authoritative:

\begin{center}
\small
\begin{tabular}{p{0.27\textwidth} p{0.29\textwidth} p{0.29\textwidth}}
\toprule
\textbf{Quantity} & \textbf{Report says} & \textbf{CSV says (authoritative)} \\
\midrule
\texttt{civicplus\_clean} / \texttt{civicplus\_akamai} & 27 / 1 & 26 / 2 \\
Overall tractability tier & \texttt{needs\_adapter} 70, \texttt{access\_blocked} 2 &
\texttt{needs\_adapter} 69, \texttt{access\_blocked} 3 \\
\texttt{minutes\_access} (report \S B) & \texttt{clean} 66, \texttt{unknown} 23,
\texttt{akamai\_403} 2 & \texttt{clean} 78, \texttt{unknown} 10, \texttt{akamai\_403} 3 \\
\texttt{in\_nzlud} filled & 8 pilot cities, 4 yes / 4 no & 11 rows filled, 4 yes / 7 no \\
\bottomrule
\end{tabular}
\end{center}

The third row is the report contradicting \emph{itself}: its own methods caveat already carries
the post-sweep numbers while its summary \S B was never updated. Its \S A table sums the
CivicPlus class to 28 either way.

\item \textbf{The zoning assessment's ``4 of 109 met the deadline'' is not reproducible from
\texttt{hcd.csv}.} \texttt{zoning\_envelope\_assessment.md} writes that only four Bay Area
jurisdictions met the 2023-01-31 deadline --- Alameda, Emeryville, San Francisco, San Leandro ---
and that \emph{none of the 14} did. Filtering \texttt{hcd.csv} to the 109 ABAG rows gives neither
reading. \textbf{Seven} jurisdictions have a \texttt{Date Received} strictly before the deadline
(Alameda 2022-11-16, San Leandro 2022-12-09, Sebastopol 2023-01-09, Emeryville 2023-01-20,
Fremont 2023-01-23, Berkeley 2023-01-24, Rohnert Park 2023-01-26), with San Francisco landing
exactly on 2023-01-31; and only \textbf{one} --- Alameda --- has an HCD \texttt{Reviewed Date}
before the deadline (2022-12-20). Two of the seven, Berkeley and Fremont, are in the 14-city set,
so the claim that none of the 14 adopted in time is contradicted by the data file. The
substantive point survives either way, since HCD's \emph{finding} is what ends exposure and only
Alameda had one before the deadline --- but the sentence should be corrected.

\item \textbf{The two HCD panels disagree by construction.} The 25-row panel anchors at
2023-01-31 and spans 0.0--38.4 months; the 14-row panel anchors at 2023-02-01 and spans
0.5--40.0 months, with eight localities appearing only in the smaller file (\S5).

\end{enumerate}

\subsection{Scripts that will not reproduce their own output}

\begin{enumerate}
\setcounter{enumi}{7}

\item \textbf{\texttt{build\_hcd\_panel.py} reads a \texttt{/tmp/} file that no longer exists.} It
takes its sample definition from \texttt{/tmp/preperiod\_sample.json}, a path outside the repo, so
the script is not re-runnable end-to-end. The 25-locality list should be committed alongside it
before the panel is rebuilt.

\item \textbf{\texttt{preperiod\_envelope\_probe/finalize.py} never encoded the
17$\rightarrow$15 correction.} The San Ramon and Daly City kills were applied to
\texttt{preperiod\_envelope.csv} \emph{in place}, not through the script: \texttt{finalize.py}
still sets San Ramon to \texttt{confidence = med} and Daly City to a valid text source, so
re-running it as written would regenerate the stale 17/25 file. Fold the two corrections into its
\texttt{CORR} table first. (It also reads \texttt{/tmp/preperiod\_sample.json}.) The original
report's headline ``17/25 (68\%)'' is superseded; a one-line correction note pointing to 15/25
sits at its top, and the rest of its method and platform diagnosis still stands.

\item \textbf{Three more hard-coded paths.} \texttt{archive\_depth\_probe/api\_probe.py} writes to
\texttt{/tmp/api\_probe\_results.json} and \texttt{assemble\_api.py} reads
\texttt{/tmp/api\_targets.json}, neither in the repo --- though the verified results are
transcribed literally into \texttt{assemble\_api.py}, so the merge itself is auditable. Both
CivicPlus scripts read their 14-locality target list from \texttt{/tmp/cp14.json}, and
\texttt{probe.py} additionally writes to a \texttt{/Users/danpost/\dots} path that does not exist
on this machine, so it is broken as written (\texttt{final\_pass.py}, writing to
\texttt{/tmp/cp14\_final.json}, is fine). \texttt{wrluri\_crosscheck.py} hard-codes the same
defunct \texttt{/Users/danpost/\dots} Dropbox root and needs the data root taken from the
environment.

\item \textbf{Two reports are stale in status rather than content.}
\texttt{zoning\_map\_form\_report.md} carries no supersession banner, but its numbers ---
67/28/13/1, the 55-versus-12 split, 71 apparent downloads --- all re-derive from the CSV.
\texttt{minutes\_platform\_pilot\_report.md}'s platform calls are superseded by the census
wherever they differ, and its Akamai framing is narrower than what was later measured: the
28-host sweep found only two blocked CivicPlus sites region-wide, so Fremont is an exception
rather than evidence about the platform. Fremont's own 403 is confirmed in both. The two
consolidating documents these probes originally reported into ---
\texttt{output/DATA\_STATUS.md} and \texttt{output/final\_recon\_bundle\_report.md} --- were
deleted on 2026-08-30; their content is folded into this memo.

\end{enumerate}

\subsection{What is unverified across the board}

No minutes document was retrieved anywhere in the recon except the two evidence files (the
Saratoga PDF and the Daly City text). Every depth year in \texttt{archive\_depth.csv} rests on a
listing or an API metadata record --- the 54 \texttt{low}-confidence rows and the 30
\texttt{type\_ambiguous} rows have never been checked against an actual document, and the
migration-cliff years likewise rest on feed entries. No zoning layer in
\texttt{zoning\_map\_form.csv} has been downloaded, opened, or joined to a parcel;
\texttt{download\_apparent} records that a path \emph{looked} available, not that it works. Four
localities may have no written minutes online at all --- Vallejo (agendas and video to 2006, no
minutes-type files), Sebastopol (video only), Benicia (council minutes not located) and Dixon
(1998--2009 on DVD at City Hall) --- candidates for the deferred physical-records workstream. The
remaining access-blocked or unreachable sites are Fremont, Portola Valley and San Jose (Akamai),
Redwood City and South San Francisco (unsupported platform), Fairfield and Suisun (403), and
American Canyon and Calistoga (refused connections, possibly transient). Also unverified: San
Jose's by-right read (blocked by Municode; the city posts a static Title 20 PDF that would
resolve it), the QCode paths for South San Francisco and Burlingame, every NZLUD height field for
a city whose code defers height to a General Plan map, all overlay districts, Berkeley's and the
City of San Mateo's minutes platforms, Cupertino's true minutes source, the city of Sonoma's
exact pre-2017 year, and whether Daly City genuinely seldom posts minutes or merely posts them
elsewhere.

\end{document}
